\documentclass[11pt,a4paper]{article}
\usepackage[utf8]{inputenc}
\usepackage[T1]{fontenc}
\usepackage{amsmath,amssymb}
\usepackage[margin=1in]{geometry}
\usepackage{hyperref}

\title{IOI 2024 -- Hieroglyphs}
\author{Solution}
\date{}

\begin{document}
\maketitle

\section{Problem Summary}
Given two sequences $A$ and $B$, find their universal common subsequence (UCS), or report that none exists.

\section{Solution}

\subsection{Filtering and Compression}
Elements that appear in only one sequence can never be part of a common subsequence, so we remove them first. The remaining values are coordinate-compressed.

\subsection{Candidate Construction}
The official full solution constructs a candidate subsequence by comparing how often each value still appears in the suffixes of $A$ and $B$. Whenever the next forced symbol can only come from one side, we append it and advance that sequence. If both sides can force different next symbols, no UCS exists.

\subsection{Verification}
The candidate is then verified with a structural criterion: there must be no ``single crossing'' between the two sequences with respect to the candidate. This is checked using:
\begin{itemize}
  \item earliest-occurrence index vectors;
  \item reversed sequences;
  \item a range minimum query structure.
\end{itemize}
If the candidate passes the crossing test in both directions, it is the UCS; otherwise the answer is $-1$.

\section{Complexity}
\begin{itemize}
  \item Filtering and candidate construction: near-linear.
  \item Verification with RMQ and index vectors: near-linear.
  \item Total: $O((n+m)\log(n+m))$ due to sorting / balanced-tree operations.
\end{itemize}

\end{document}
